Sirius is the brightest star of the night sky during the winter of the northern hemisphere and the summer
of the southern hemisphere. It has an apparent visual magnitude $m_V=-1.47$ mag. Sirius B is the binary
companion of Sirius A and its magnitude is $m_V=8.5$ mag. Their orbital period is 50 years and the parallax
of the system is 350 milliarcseconds, while their angular separation is 7 arcseconds.
You can assume that Sirius A has a surface temperature of $T=10{,}000$ K. The temperature of Sirius B is
estimated to be 25,000 K. Both stars have a bolometric correction equal to $-0.4$. An absolute bolometric
magnitude of zero corresponds to a luminosity of $3\times10^{28}$ W. The orbit is assumed to be circular,
and the line of sight is perpendicular to the orbital plane of the system.
\begin{description}
\item[(A)] Calculate the luminosity and radius of Sirius A and B. \hfill\mbox{[12 points]}
\item[(B)] What is the orbital separation of the system in meters? \hfill\mbox{[8 points]}
\item[(C)] What is the total mass of the system? \hfill\mbox{[8 points]}
\item[(D)] Assuming that the mass of Sirius A is twice the mass of Sirius B, what is the mean density of Sirius B? \hfill\mbox{[6 points]}
\item[(E)] In the H-R diagram shown in question 10, which star number (1 / 2 / 3) will be the closest to the star Sirius B? \hfill\mbox{[6 points]}
\end{description}
\ioaafig{2023_TH_Q10-f1.pdf}
\textit{Data from the accompanying constants sheet.}
$\sigma=5.670\times10^{-8}\,\mathrm{W\,m^{-2}\,K^{-4}}$, $1\,\mathrm{pc}=3.086\times10^{16}\,\mathrm{m}$,
$1\,\mathrm{au}=1.496\times10^{11}\,\mathrm{m}$, and $M_\odot=1.988\times10^{30}\,\mathrm{kg}$.
